\documentclass[11pt]{article}

\usepackage[margin=1in]{geometry}
\usepackage{amsmath,amssymb,amsthm}
\usepackage{mathtools}

\newtheorem{theorem}{Theorem}
\newtheorem{lemma}{Lemma}
\newtheorem{remark}{Remark}

\DeclareMathOperator*{\supremum}{sup}

\begin{document}

\title{m=3, two jobs zero}
\author{}
\date{}
\maketitle

\section{Problem Introduction}

Let
\[
0\le \alpha\le \frac12 \quad\text{and}\quad
a,b,c,d,e,f,p>0.
\]
Set
\[
S:=a^2+b^2+c^2+d^2+e^2+f^2+p^2.
\]
Consider the function
\[
\lambda(a,b,c,d,e,f,p)(\alpha)
=
\frac{
(a+d)^2+(b+e)^2+(c+f+p)^2+(2\alpha-1)S
}{
(a+b+c)^2+(d+e+f)^2+p^2+(2\alpha-1)S
}
\]
subject to
\[
\begin{aligned}
a+b+c&\ge d+e,\\
d+e+f&\ge a+b,\\
p&\ge a+b,\\
p&\ge d+e,\\
a+d&\ge c+f,\\
b+e&\ge c+f,\\
a+d&\ge b,\\
b+e&\ge a,\\
c+f+p&\ge a,\\
c+f+p&\ge b.
\end{aligned}
\]


\section{First Simplification}

We normalize the problem by setting
\[p=1.
\]

The final two constraints are redundant. Indeed, since
\[
p\ge a+b
\]
and \(a,b>0\), we have \(p\ge a\) and \(p\ge b\). Hence
\[
c+f+p\ge a,\qquad c+f+p\ge b
\]
hold "automatically".

\section{Expansion of the Ratio}

Define
\[
L:=ab+ac+bc+de+df+ef
\quad \text{and}\quad
Q:=ad+be+cf+p(c+f).
\]
Expanding the denominator of $\lambda$ gives
\[
\begin{aligned}
& (a+b+c)^2+(d+e+f)^2+p^2+(2\alpha-1)S\\
&=
2\alpha S+2(ab+ac+bc+de+df+ef)\\
&=
2(\alpha S+L).
\end{aligned}
\]
Similarly for the numerator,
\[
\begin{aligned}
& (a+d)^2+(b+e)^2+(c+f+p)^2+(2\alpha-1)S\\
&=
2\alpha S+2(ad+be+cf+p(c+f))\\
&=
2(\alpha S+Q).
\end{aligned}
\]
Thus
\[
\lambda(a,b,c,d,e,f,p)(\alpha)
=
\frac{\alpha S+Q}{\alpha S+L}
=
1+\frac{Q-L}{\alpha S+L}.
\]
Therefore the problem is equivalent to maximizing
\[
\frac{Q-L}{\alpha S+L}.
\]

Only one term in \(Q\) involves \(p\), namely
\[
p(c+f).
\]
The denominator contains the compensating term
\[
\alpha p^2.
\]
Thus, making \(p\) large cannot make the ratio
arbitrarily large.

Let
\[
x=c+f.
\]
The constraints
\[
a+d\ge x,\qquad b+e\ge x
\]
say that the two sums \(a+d\) and \(b+e\) must each carry at least the
mass \(x\).  We do not set these inequalities tight at this point.
Instead, the next section derives the active boundary algebraically and
only then checks that any slack in these two inequalities lowers the
ratio on that reduced face.
The same applies to the constraint \(b+e\ge a\): it is not assumed
active here.  After \(a+d=b+e=x\) and \(c=f=x/2\) have been reduced to,
the remaining dependence is through \(y=a+b\), and the numerator
\[
px-\frac74x^2+2xy-y^2
\]
is maximized at \(y=x\); the denominator then forces \(ab=0\), so on
the symmetric branch \(a=x,b=0,e=x\), giving \(b+e=a\).

\section{Constraint Reduction}

We now record the algebraic reduction leading to the active boundary.

\begin{lemma}[Reduction to the active boundary]
For \(0<\alpha\le \frac12\),
\[
\frac{Q-L}{\alpha S+L}
\le
\sup_{t>0}
\frac{t-\frac34t^2}{\alpha+\left(1+\frac52\alpha\right)t^2}.
\]
\end{lemma}

\begin{proof}
Introduce the sums
\[
u=a+d,\qquad v=b+e,\qquad x=c+f.
\]
The constraints include
\[
u\ge x,\qquad v\ge x.
\]
Also put
\[
m=a+b,\qquad n=d+e.
\]
The constraints
\[
p\ge a+b,\qquad p\ge d+e
\]
become
\[
p\ge m,\qquad p\ge n.
\]
Finally,
\[
a+d\ge b,\qquad b+e\ge a
\]
become
\[
u\ge b,\qquad v\ge a.
\]

For fixed \(p\) and \(x=c+f\), expanding \(Q-L\) in the variables
\(u,v,m,c\) gives
\[
\begin{aligned}
Q-L
={}&px-c^2+c(u+v+x-2m)-m^2\\
&+m(u+v+x)-uv-x(u+v).
\end{aligned}
\]
As we set
\[
a+d=x,\qquad b+e=x.
\]
we can analyze the remaining constraint \(b+e\ge a\). Choose the maximizing "split"
\[
c=f=\frac x2,
\]
which is the maximal feasible choice for fixed \(x\), since it maximizes \(cf\) (numerator)
and minimizes \(c^2+f^2\) (denominator).  Then
\[
d=x-a,\qquad e=x-b,
\]
put
\[
y=a+b.
\]
Then inserting gives
\[
Q-L
=
px-\frac74x^2+2xy-y^2.
\]
For fixed \(x\) and \(p\), this quadratic in \(y\) is maximized at
\[
y=x.
\]
(as can be seen by differentiation).

Thus the maximizing satisfies
\[
a+b=x.
\]
Here, the denominator is
\[
\alpha S+L
=
\alpha p^2+\left(1+\frac52\alpha\right)x^2
+(2-4\alpha)ab.
\]
Since \(0<\alpha\le \frac12\), the coefficient \(2-4\alpha\) is
nonnegative.  Hence the denominator is minimized by
\[
ab=0.
\]
Consequently either
\[
(a,b)=(x,0)
\qquad\text{or}\qquad
(a,b)=(0,x).
\]
The two cases are symmetric.  We choose the first one, which implies:
\[
e=x,\qquad b+e=x=a.
\]
Thus, on this symmetric branch, the constraint \(b+e\ge a\) is tight.

Finally, we need to check that this does not benefit from slack in
\(u\ge x\) or \(v\ge x\). We write
\[
u=x+r,\qquad v=x+s,
\qquad r,s\ge0.
\]
this gives
\[
(a,b,c,d,e,f,p)
=
\left(x,s,\frac x2,r,x,\frac x2,p\right),
\]
so that
\[
a+d=x+r,\qquad b+e=x+s,\qquad c+f=x.
\]
For this family the excess numerator is
\[
\begin{aligned}
Q-L
&=px-\frac34x^2-\frac{x}{2}(r+s),
\end{aligned}
\]
whereas the denominator is
\[
\begin{aligned}
\alpha S+L
&=
\alpha p^2+\left(1+\frac52\alpha\right)x^2
+\frac{3x}{2}(r+s)+\alpha(r^2+s^2).
\end{aligned}
\]
Thus every positive slack \(r\) or \(s\) subtracts from the numerator
and adds to the denominator.  Consequently, for fixed \(x\) and \(p\),
the largest value in this slack family is attained at
$
r=s=0.
$

Substituting the active equalities gives
\[
(a,b,c,d,e,f,p)
=
\left(x,0,\frac x2,0,x,\frac x2,p\right).
\]
We normalize p to get
\[
\left(t,0,\frac t2,0,t,\frac t2,1\right).
\]
On this boundary,
\[
Q-L=t-\frac34t^2
\]
and
\[
\alpha S+L
=
\alpha+\left(1+\frac52\alpha\right)t^2.
\]
Therefore every feasible point is bounded above by the one-variable
problem
\[
\sup_{t>0}
\frac{t-\frac34t^2}{\alpha+\left(1+\frac52\alpha\right)t^2}.
\]
\end{proof}

\section{Max t}

Define
\[
f_\alpha(t)
=
\frac{t-\frac34t^2}{\alpha+\left(1+\frac52\alpha\right)t^2},
\qquad t>0.
\]
Set
\[
\beta=1+\frac52\alpha.
\]
Then
\[
f_\alpha(t)=\frac{t-\frac34t^2}{\alpha+\beta t^2}.
\]
Differentiate:
\[
f_\alpha'(t)
=
\frac{
\left(1-\frac32t\right)(\alpha+\beta t^2)
-
\left(t-\frac34t^2\right)(2\beta t)
}{
(\alpha+\beta t^2)^2
}.
\]
The numerator simplifies to
\[
\alpha-\frac32\alpha t-\beta t^2.
\]
Thus the critical point satisfies
\[
\beta t^2+\frac32\alpha t-\alpha=0.
\]
The positive root is
\[
t_\alpha
=
\frac{
-\frac32\alpha+\sqrt{\frac94\alpha^2+4\alpha\beta}
}{
2\beta
}.
\]
Since
\[
\beta=1+\frac52\alpha,
\]
we have
\[
\frac94\alpha^2+4\alpha\beta
=
4\alpha+\frac{49}{4}\alpha^2
=
\frac14\alpha(16+49\alpha).
\]
Therefore
\[
t_\alpha
=
\frac{\sqrt{\alpha(16+49\alpha)}-3\alpha}{4+10\alpha}.
\]

So the supremum is

\[
\lambda_\alpha(t_\alpha)
=
\frac{2}{\sqrt{\alpha(16+49\alpha)}+3\alpha}.
\]


\section{Appsroaching family}

The maximizing boundary point has zero coordinates:
\[
\left(
t_\alpha,0,\frac{t_\alpha}{2},0,t_\alpha,\frac{t_\alpha}{2},1
\right).
\]
Which is approached by the family of points
\[
\left(
t_\alpha,\varepsilon,\frac{t_\alpha}{2},
\varepsilon,t_\alpha,\frac{t_\alpha}{2},1
\right)
\]
as \(\varepsilon\to0\).  This family is feasible for every \(\varepsilon>0\).

\end{document}
